\documentclass[11pt]{article}
\usepackage[a4paper,margin=26mm]{geometry}
\usepackage{amsmath,amssymb,amsthm,hyperref}
\newtheorem{theorem}{Theorem}
\newtheorem{lemma}[theorem]{Lemma}
\newtheorem{corollary}[theorem]{Corollary}
\newcommand{\PP}{\mathcal P}
\title{A permutation-fair collection with one polynomial-size die}
\author{Research notes, 30 September 2026}
\date{}
\begin{document}
\maketitle
\noindent\textbf{Status.} Complete proof with independent mathematical reviews completed
by three collaborating agents. These reviews are not machine-checked formal
verification. The theorem disproves the manuscript's conjecture that every
die in a permutation-fair collection must have exponentially many faces.
The other dice in the construction can be arbitrarily large; this does not
improve the optimum total size.

\begin{theorem}\label{thm:main}
There is an absolute constant $C$ such that, for every $n\ge2$, there
is a permutation-fair collection of $n$ finite uniform dice, with pairwise
disjoint face labels, in which one distinguished die has at most $Cn^2$
faces. The other $n-1$ dice can all be chosen with the same number of faces.
\end{theorem}

The argument first builds positive rational piecewise-constant densities
for the other dice and a finite uniform distinguished die. It then replaces
each constant-density interval by a finite permutation-fair block.

Write $\PP_d$ for the real polynomials of degree at most $d$.

\begin{lemma}[Regular real starting rule]\label{lem:real}
For every $m\ge1$ there is a real equal-weight uniform-interval quadrature
of degree $m$ with
\[
 m^2\le K=O(m^2),\qquad 0<u_1<\cdots<u_K<1.
\]
\end{lemma}
\begin{proof}
The classical Chebyshev-type quadrature bound gives an equal-weight
degree-$(2m+2)$ rule with $K_0=O(m^2)$ nodes in $[0,1]$.
For example, use the uniform-weight case of Gilboa--Peled, Theorem 1.1:
\url{https://arxiv.org/abs/1507.01505}.
It has at least $m+2$ distinct support points: otherwise the square of
its support-vanishing polynomial would have degree at most $2m+2$,
zero quadrature sum, and positive integral. At least $m$ support points
are interior. Repeat the rule $\lceil m^2/K_0\rceil$ times so its
cardinality $K$ is at least $m^2$ and still $O(m^2)$.

Select $m$ coordinates at distinct interior support points as pivots.
The Jacobian of the first $m$ moment sums with respect to them is an
invertible scaled Vandermonde matrix. The implicit function theorem permits
arbitrary small perturbations of the remaining coordinates, compensating
with the pivots to preserve these moment sums. The pivots remain interior
and mutually distinct. Move every free endpoint slightly inward.

Free--free collisions are proper affine hyperplanes. If a free coordinate
and a pivot coincide initially, varying just that free coordinate changes
its matching pivot with derivative $-1$, since their Jacobian columns
coincide; all other pivot derivatives are zero. Their difference therefore
has derivative $-2$, so that collision is a proper analytic zero set.
Initially absent collisions stay absent in a sufficiently small
neighborhood. A finite union of proper analytic zero sets cannot cover
the open set of permitted inward perturbations. Choose a collision-free
perturbation, then sort the nodes.
\end{proof}

\begin{lemma}[Rational separated moment bumps]\label{lem:bumps}
For the nodes in Lemma~\ref{lem:real}, one can choose pairwise disjoint open
intervals $I_q\subset(0,1)$ around $u_q$, smaller neighborhoods
$V_q\subset I_q$ around $u_q$, and bounded rational piecewise-constant
functions $h_q$, with rational breakpoints and support in $I_q$, such that
for every $t\in V_q$ and every $g\in\PP_{m-1}$,
\[
 \int_0^t h_q(x)g(x)\,dx=g(0),\qquad
 \int_t^1 h_q(x)g(x)\,dx=-g(0).
\]
In particular $\int_0^1h_q(x)g(x)\,dx=0$.
\end{lemma}
\begin{proof}
Choose disjoint rational-endpoint neighborhoods $I_q$. Inside each $I_q$,
choose $m$ disjoint rational closed subintervals strictly to the left of
$u_q$, and $m$ strictly to its right, all of positive length. Their union
misses a neighborhood $V_q$ of $u_q$.

For any $m$ ordered disjoint intervals $J_1<\cdots<J_m$, the rational
matrix
\[
 M_{\ell a}=\int_{J_a}x^\ell\,dx\qquad
 (0\le\ell\le m-1,\ 1\le a\le m)
\]
is invertible. Indeed, multilinearity of the determinant gives
\[
 \det M=\int_{J_1\times\cdots\times J_m}
         \prod_{a<b}(x_b-x_a)\,dx_1\cdots dx_m>0.
\]
Solve $M\gamma=(1,0,\ldots,0)^T$ for the left intervals and
$M\gamma=(-1,0,\ldots,0)^T$ for the right intervals. Use these rational
coefficients as the constant values of $h_q$ on those intervals, and zero
elsewhere. The claimed identities follow by linearity. No nonnegativity
is claimed for $h_q$ itself.
\end{proof}

\begin{lemma}[Order response of disjoint moment bumps]\label{lem:response}
Choose $t_q\in V_q$. Let $Q_1,\ldots,Q_m$ be pairwise disjoint subsets
of $[K]$. Suppose
\[
 f_j(x)=1+\sum_{q\in Q_j}\delta_{j,q}h_q(x),\qquad j\in[m],
\]
are nonnegative densities. Let $X_j$ have density $f_j$, independently,
and let $X_0$ be independent and uniform on $t_1,\ldots,t_K$.
Define functionals on $\PP_{m-1}$ by
\[
 B_j(a)=\frac1K\sum_{q\in Q_j}\delta_{j,q}a(t_q).
\]
If all $B_j$ agree with one functional $B$, then every full ordering
placing $X_0$ after exactly $k$ of the old variables has probability
\begin{equation}\label{eq:response}
 \frac1K\sum_{q=1}^Kp_k(t_q)+B(p_k'),\qquad
 p_k(t)=\frac{t^k(1-t)^{m-k}}{k!(m-k)!}.
\end{equation}
\end{lemma}
\begin{proof}
Fix a full order and condition on $X_0=t_\ell$. Expand the product of
the old densities. The term choosing no bumps has probability
$p_k(t_\ell)$.

Consider a term choosing bumps for a nonempty set of old variables,
and expand each chosen density's sum over its indices. All chosen bumps
have distinct indices because the $Q_j$ are disjoint. Their supporting
intervals are therefore disjoint and ordered. If their order is inconsistent
with the fixed full order, the term is zero. Otherwise their mutual ordering
is automatic throughout those intervals. Integrate out the old variables
whose density contribution was the constant $1$. If there are $b$ chosen
bumps, the resulting kernel is a product of powers of the gaps between
the chosen coordinates, $0$, $1$, and $t_\ell$, divided by factorials.
It is a polynomial of total degree $m-b$ in the chosen coordinates and
$t_\ell$ on each of the relevant regions.

At most one supporting interval contains $t_\ell$, namely $I_\ell$.
If $b\ge2$, at least one chosen bump has its entire support on one side
of $t_\ell$. For fixed remaining coordinates, its integration runs over
its entire support against a polynomial of degree at most $m-b\le m-1$.
The full moment identities of Lemma~\ref{lem:bumps} make this integral
zero. The same argument eliminates a one-bump term whose index is not
$\ell$. Thus only the one-bump terms cut by the distinguished face remain.

We calculate them. If the perturbed old variable is before $X_0$, let
$a$ old variables precede it and $b$ lie between it and $X_0$, so
$a+b=k-1$. Its kernel is
\[
 g^-_{a,b}(x,t)=
 \frac{x^a(t-x)^b(1-t)^{m-k}}{a!b!(m-k)!}.
\]
Its integration against the left part of $h_q$ is $g^-_{a,b}(0,t)$.
If the perturbed variable is after $X_0$, let $a$ old variables lie
between them and $b$ follow it, so $a+b=m-k-1$. Its kernel is
\[
 g^+_{a,b}(x,t)=
 \frac{t^k(x-t)^a(1-x)^b}{k!a!b!},
\]
whose integral against the right part of $h_q$ is $-g^+_{a,b}(0,t)$.
These evaluated kernels have degree at most $m-1$ in $t$.

Average over $X_0$ and sum over the old variables. Since all $B_j=B$,
the resulting correction is $B$ applied to the sum of these kernel
polynomials. The binomial theorem gives
\begin{align*}
 \sum_{a+b=k-1}g^-_{a,b}(x,t)
  &=\frac{t^{k-1}(1-t)^{m-k}}{(k-1)!(m-k)!},\\
 \sum_{a+b=m-k-1}g^+_{a,b}(x,t)
  &=\frac{t^k(1-t)^{m-k-1}}{k!(m-k-1)!}.
\end{align*}
Use the first expression only when $k\ge1$ and the second only when
$k\le m-1$. Their difference is $p_k'(t)$, including the endpoint
cases. This proves~\eqref{eq:response}.
\end{proof}

\begin{proof}[Proof of Theorem~\ref{thm:main}: continuous intermediate model]
Put $m=n-1$ and choose the regular real rule and bumps above. Since
$K\ge m^2$, choose disjoint sets $Q_1,\ldots,Q_m$, each of size $m$.
Choose rational $t_q\in V_q$ arbitrarily close to $u_q$. Their ordering
remains strict. For $1\le s\le m$, put
\[
 \mu_s(t)=\frac1K\sum_{q=1}^Kt_q^s,\qquad
 a_s(t)=\frac{1/(s+1)-\mu_s(t)}s.
\]
For each $j$, solve the rational Vandermonde system
\[
 \frac1K\sum_{q\in Q_j}\delta_{j,q}t_q^{s-1}=a_s(t)
       \qquad(1\le s\le m).
\]
The matrix is invertible because the $t_q$ are distinct. The solution
depends continuously on $t$ near $u$ and tends to zero as $t\to u$,
because $a_s(u)=0$. All $h_q$ are fixed bounded functions. Taking the
rational approximation sufficiently close therefore guarantees
$f_j(x)\ge1/2$ everywhere, where
\[
 f_j(x)=1+\sum_{q\in Q_j}\delta_{j,q}h_q(x).
\]
Their integrals are $1$ by the full zeroth-moment identity for the bumps.
They are thus positive rational piecewise-constant probability densities.

The functionals $B_j$ agree on the monomial basis of $\PP_{m-1}$,
and hence agree with one functional $B$. For every $p\in\PP_m$,
\[
 B(p')=\int_0^1p(t)\,dt-\frac1K\sum_{q=1}^Kp(t_q),
\]
as follows by checking $p(t)=t^s$, $1\le s\le m$, and constants.
Apply Lemma~\ref{lem:response} to $p_k$. Every full order has probability
\[
 \int_0^1p_k(t)\,dt=\frac1{(m+1)!}.
\]
This proves permutation fairness for the continuous old variables and the
$K$-face distinguished variable.
\end{proof}

\begin{proof}[Finite realization, completing the theorem]
Partition $[0,1]$ at every rational breakpoint of every $f_j$ and at all
$t_q$. Ignore singleton boundary points, which have probability zero for
the old variables. On each remaining interval $J_g$, each $f_j$ is a
positive rational constant. The masses
\[
 \alpha_{j,g}=\int_{J_g}f_j(x)\,dx
\]
are positive rationals, and $\sum_g\alpha_{j,g}=1$ for each $j$.
Choose a common positive integer $D$ clearing all their denominators.

Take any finite equal-sided permutation-fair collection of $m$ dice,
with $F$ faces each, whose existence is already proved in the manuscript.
In interval $J_g$, put a copy of its fair word in which letter $j$ is
blown up by the positive integer $D\alpha_{j,g}$. This block is
permutation-fair, and contains $FD\alpha_{j,g}$ faces of old die $j$.
Concatenate the blocks in interval order, with one distinguished face
inserted at each boundary $t_q$. Every old die has exactly $FD$ faces;
the distinguished die has exactly $K$.

For each old die, its interval probabilities equal $\alpha_{j,g}$,
exactly as in the continuous model. Conditional on the selected intervals,
old variables in distinct intervals have a determined relative order.
Those in a common interval are uniformly ordered in the continuous model,
because their conditional distributions are all uniform on that interval.
They are also uniformly ordered in the discrete model, by permutation
fairness of that block and its restrictions. Different interval groups
involve disjoint sets of independent dice, so their internal orders combine
independently. Thus the entire order law, including the distinguished
faces at the interval boundaries, is unchanged. Assign successive distinct
integer labels to the concatenated word to obtain the required finite dice.
\end{proof}

\paragraph{What is and is not proved.}
The distinguished die has $O(n^2)$ faces, while the sizes of the other
dice are not bounded quantitatively here. The construction uses arbitrarily
accurate rational approximation and correspondingly large denominator
clearing only for those other dice. It does not conflict with the manuscript's
result that almost all dice must be exponentially large, nor with the new
exponential obstruction for rational scalar equal-weight interval designs.
The local CDF vectors are deliberately different for different old dice.

\end{document}
